\documentclass[12pt]{article}

\usepackage[margin=1.2in]{geometry}
\usepackage{amsmath, amssymb}
\usepackage{hyperref}
\usepackage{booktabs}
\usepackage{array}
\usepackage{tcolorbox}
\usepackage{parskip}
\usepackage[T1]{fontenc}
\usepackage{xcolor}
\usepackage{authblk}

\tcbuselibrary{skins, breakable}

\definecolor{navyblue}{RGB}{10,30,100}
\definecolor{crfgreen}{RGB}{0,100,60}
\definecolor{softgreen}{RGB}{180,230,210}
\definecolor{physgold}{RGB}{140,100,0}
\definecolor{softgold}{RGB}{255,240,180}
\definecolor{loopred}{RGB}{160,30,30}
\definecolor{softred}{RGB}{255,210,210}
\definecolor{midblue}{RGB}{30,80,160}
\definecolor{softblue}{RGB}{180,210,240}
\definecolor{mutedgray}{RGB}{90,90,90}

\hypersetup{
  colorlinks=true,
  linkcolor=navyblue,
  urlcolor=midblue,
  citecolor=navyblue
}

\newtcolorbox{predbox}{
  colback=softgold!60, colframe=physgold!80,
  fonttitle=\bfseries, title={Prediction P0},
  breakable
}

\newtcolorbox{derivedbox}[1][]{
  colback=softgreen!50, colframe=crfgreen!60,
  fonttitle=\bfseries\color{crfgreen!20!black},
  title={#1}, breakable
}

\newtcolorbox{resultbox}[1][]{
  colback=softblue!40, colframe=midblue!60,
  fonttitle=\bfseries\color{navyblue},
  title={#1}, breakable
}

\begin{document}

%% ── TITLE ──────────────────────────────────────────────────────────────
\begin{center}
{\LARGE\bfseries\color{navyblue}
Constrained Reality Framework:}\\[0.5em]
{\Large\bfseries\color{navyblue}
Deriving Quantum Measurement, Gravity, and\\
the Gravitational Constant from a Single Primitive}\\[0.7em]
{\normalsize Ougit Jarittum}\\[0.3em]
{\small\color{mutedgray} March 2026 \quad
Zenodo: \href{https://doi.org/10.5281/zenodo.18977059}
{10.5281/zenodo.18977059}}\\[0.4em]
{\small\color{mutedgray} Version 3 ---
Born Rule as theorem, Schwarzschild derived,
$G = \ln(2)\cdot c^2/D_0$.}
\end{center}

\bigskip

%% ── ABSTRACT ────────────────────────────────────────────────────────────
\noindent\textbf{Abstract.}
The Constrained Reality Framework (CRF) proposes that reality is the
continuous outcome of a single process: the probabilistic resolution of
latent informational possibility under constraint. It derives the
entirety of its structure from one primitive --- Distinction ($\delta$)
--- without additional axioms.

This paper reports three results beyond the original framework.
\textit{First}: the Born Rule is established as a theorem of CRF via
Gleason's Theorem applied to the $\delta$-frequency measure on
Possibility Space $P$ (dimension $\geq 3$ from SO(3) symmetry). Born
Rule probabilities are not imported --- they are the unique statistics
of Resolution at minimal constraint.
\textit{Second}: the Schwarzschild metric is derived from the
non-linear $\mathcal{R}$-feedback equation
$d\mathcal{C}/dr = -(k/\mathcal{C})\,e^{-\mathcal{C}^2/\alpha^2 D_0}$,
which follows from $\theta = 1 - e^{-D_{\mathrm{KL}}/D_0}$ and the
quadratic scaling $D_{\mathrm{KL}} \propto \mathcal{C}^2$ near the
gravitational horizon.
\textit{Third}: Newton's gravitational constant is derived as
$G = \ln(2)\cdot c^2/D_0$, where $D_0 = \mathrm{Scale}(P)$ is the
information-scale of the vacuum. No free parameter of gravity remains.

A testable prediction (P0) follows from the framework: a measuring
device with greater measurement history causes faster wave function
collapse. Formally, $\tau_1 < \tau_2$ when $N_{\mathrm{eff},1} >
N_{\mathrm{eff},2}$, where $\tau$ is collapse time and
$N_{\mathrm{eff}} = \sum_i \Delta D_{\mathrm{KL},i}$. The formula
$\tau = \tau_0 e^{-\kappa N_{\mathrm{eff}}}$ with
$\kappa = k_B\ln(2)/H_{\mathrm{material}}$ contains no free
parameters. An experimental protocol using SNSPD detectors is
proposed.

\bigskip
\tableofcontents
\newpage

%% ── 1 ───────────────────────────────────────────────────────────────────
\section{The Question and the Primitive}

This framework began with an observation that preceded any formalism:
every stable thing we encounter is stable \textit{under constraint}.
A rock, a cell, a mind, a culture --- each is stable because something
outside itself holds it in place. None of these stabilities are free.

The question is: is there a ground condition for stability that does
not depend on anything outside itself?

Following this question inward leads to a more primitive question:
what has to be true for anything to be different from anything else?
If nothing is different, there are no states to be stable or unstable.
Distinguishability is the precondition of everything we can ask about.

The answer requires only one undefined term: Distinction ($\delta$),
the condition that makes any state distinguishable from any other.
$\delta$ has no ground outside itself. It is the self-grounding
primitive of CRF --- declared, as Euclid declared his points and
lines, as the irreducible starting condition.

From $\delta$, without adding further assumptions, the entire
structure follows.

%% ── 2 ───────────────────────────────────────────────────────────────────
\section{The Derivation Chain}

\subsection{From $\delta$ to Probability}

$\delta$ operating on the undifferentiated produces $P$ ---
Possibility Space, the accumulated trace of all $\delta$-activity:
$P = \delta(\varnothing) = \{s \mid s\text{ is distinguished but
unresolved}\}$.

The relative frequency with which $\delta$ distinguishes state $s$
is probability:
\[
  \Pr(s) = \lim_{n\to\infty}
  \frac{\#\{\delta\text{-operations distinguishing }s\}}{n}
\]
Probability is not assumed. It is derived from $\delta$-frequency.

\subsection{From Probability to Resolution and Constraint}

When $\Pr(s|\mathcal{C}_{\mathrm{internal}}) \geq \theta$, state $s$
becomes actual. This is Resolution: $\mathcal{R}$. Time is the count
of Resolutions: $t = |\mathcal{R}_{\mathrm{events}}|$.

The accumulated record of all prior $\mathcal{R}$-events is
Constraint: $\mathcal{C}$. Since $\mathcal{C}$ operates entirely
through $\Pr(s|\mathcal{C})$, they are identical:
$\mathcal{C} \equiv \Pr(s|\mathcal{C})$. This identification removes
a hidden assumption that they were separate.

\subsection{Mind as Boundary and the Threshold}

Mind ($\mathcal{M}$) is a system whose $\mathcal{C}$-structure is
distinct from its surroundings:
$\mathcal{C}_{\mathrm{int}} \neq \mathcal{C}_{\mathrm{ext}}$. The
boundary between these distributions is measured by KL divergence,
which emerges from $\delta$-frequency as the $\delta$-activity gap
across $\mathcal{M}$'s boundary. The threshold is not fixed --- it
is the current state of the boundary:
\[
  \theta = 1 - e^{-D_{\mathrm{KL}}/D_0}
\]
where $D_0$ is the characteristic $D_{\mathrm{KL}}$ scale of the
system. The derivation chain is complete:
\[
  \delta \;\xrightarrow{\Pr(s)}\; \Pr(s|\mathcal{C})
  \;\xrightarrow{D_{\mathrm{KL}}}\;
  \theta = 1 - e^{-D_{\mathrm{KL}}/D_0}
  \;\xrightarrow{\mathcal{R}}\;
  t = |\mathcal{R}_{\mathrm{events}}|
\]

%% ── 3 ───────────────────────────────────────────────────────────────────
\section{Born Rule as Theorem}

\subsection{The Bridge via Gleason's Theorem}

CRF derives $\Pr(s)$ as $\delta$-frequency. Standard quantum
mechanics postulates the Born Rule $\Pr = |\langle s|\psi\rangle|^2$.
These are the same, via Gleason's Theorem (1957).

\begin{derivedbox}[Theorem: Born Rule from $\delta$-Statistics]
\textbf{Premises}: (i) Axiom 0: $\delta$-operations constitute a
counting measure on $P$. (ii) Session 5 (SO(3)): $P$ has
dimension $\geq 3$.

\textbf{Gleason's Theorem}: in a Hilbert space of dimension
$\geq 3$, every frame function (non-negative, normalized measure
on closed subspaces) has the unique form
$\mu(s) = \mathrm{Tr}(\rho\,|s\rangle\langle s|)$.

\textbf{Application}: at $\mathcal{C} = \mathcal{C}_{\min}$
(no measurement history), $\rho = |\psi\rangle\langle\psi|$ and:
\[
  \Pr(s) = |\langle s|\psi\rangle|^2
\]
This is the Born Rule --- not imported, but the unique
$\delta$-frequency measure consistent with $P$'s Hilbert
structure at $\mathcal{C}_{\min}$.

For qubits ($n=2$): the continuity of
$\theta = 1 - e^{-D_{\mathrm{KL}}/D_0}$ enforces the same
structure via the extended Gleason theorem (Busch 2003).
\end{derivedbox}

\subsection{Consequence for the Measurement Problem}

The complete conditional probability structure is:
\[
  \Pr(s|\mathcal{C}) =
  \begin{cases}
    |\langle s|\psi\rangle|^2 & \mathcal{C} = \mathcal{C}_{\min} \\
    f(\theta,\, N_{\mathrm{eff}}) & \mathcal{C} > \mathcal{C}_{\min}
  \end{cases}
\]
An observer is any system with sufficient accumulated $\mathcal{C}$
to lower the local $\mathcal{R}$-threshold. Consciousness is not
required. What causes $\mathcal{R}$ is $\mathcal{C}$, not
consciousness.

%% ── 4 ───────────────────────────────────────────────────────────────────
\section{Schwarzschild Metric and the Gravitational Constant}

\subsection{Newtonian Limit (Prior Result)}

From the $\mathcal{R}$-feedback loop (Axiom 2): $\mathcal{C}$ grows
in the direction of decreasing $\mathcal{R}$-density. In the linear
regime ($\rho_\mathcal{R} \propto 1/\mathcal{C}$):
\[
  \mathcal{C}\frac{d\mathcal{C}}{dr} = -k
  \quad\Rightarrow\quad
  \mathcal{C}(r) \approx \alpha\!\left(1 + \frac{GM}{c^2 r}\right)
  \quad (r \gg r_s)
\]
This is Newtonian gravity, derived from $\delta$ without importing
the gravitational framework.

\subsection{Schwarzschild: Non-Linear Feedback}

Near the gravitational horizon, $\mathcal{C} \gg \alpha$. The full
$\theta$-form must be used. Since $D_{\mathrm{KL}} \propto
\mathcal{C}^2/\alpha^2$ (quadratic scaling of information distance
between a peaked $\mathcal{C}$ and flat $P_{\mathrm{vac}}$):
\[
  \rho_\mathcal{R} \propto (1-\theta) = e^{-D_{\mathrm{KL}}/D_0}
  \propto e^{-\mathcal{C}^2/\alpha^2 D_0}
\]

\begin{derivedbox}[Schwarzschild Form from Non-Linear Feedback]
The non-linear $\mathcal{R}$-feedback equation:
\[
  \frac{d\mathcal{C}}{dr} = -\frac{k}{\mathcal{C}}\,
  e^{-\mathcal{C}^2/\alpha^2 D_0}
\]
integrates near $r_s$ to:
\[
  \mathcal{C}(r) = \frac{\alpha}{1 - r_s/r}
\]
The Schwarzschild line element follows from
$d\tau/dt = 1/\sqrt{\mathcal{C}/\alpha}$:
\[
  ds^2 = -\!\left(1-\frac{r_s}{r}\right)c^2 dt^2
  + \left(1-\frac{r_s}{r}\right)^{\!-1}dr^2
  + r^2 d\Omega^2
\]
All boundary conditions (C1--C4) satisfied. The horizon
$r = r_s$ is where $\mathcal{R} \to 0$: the radius at which
Resolution ceases. The universe stops deciding.
\end{derivedbox}

\subsection{The Gravitational Constant Derived}

$k$ is the rate of $\mathcal{C}$-saturation per unit $r$ at the
information horizon. Both $\alpha = \ln(2)$ (Session 8:
$\delta$-bit equivalence) and $D_0$ are already present in CRF.
The natural identification is $k = \alpha^2/D_0$:

\begin{derivedbox}[$G$ from $\delta$]
From $k = \alpha^2/D_0$ and the matching condition
$k/\alpha = GM/c^2$:
\[
  \boxed{G = \frac{\ln(2) \cdot c^2}{D_0}}
\]
$G$ is the conversion ratio between Bit-Density and Spatial
Curvature. $D_0 = \mathrm{Scale}(P)$ is the
information-scale of the vacuum --- the one measurable property
of $P$ that sets the scale of all gravitational phenomena.

No external parameter of gravity remains.

Consequence: $G(t) \propto 1/P(t)$ --- $G$ decreases as $P$
grows with $\delta$-activity. A prediction: $G$ and $\Lambda$
co-vary as $1/P(t)$, consistent with DESI 2024 observations of
$\Lambda$ evolution.
\end{derivedbox}

%% ── 5 ───────────────────────────────────────────────────────────────────
\section{The Prediction (P0)}

\subsection{Formula: Zero Free Parameters}

A measuring device $D$ is a system with internal $\mathcal{C}_D$
maintained distinctly from its environment. Each measurement adds
to $\mathcal{C}_D$. Two devices $D_1$, $D_2$ with histories
$N_1 > N_2$ have $\mathcal{C}_{D_1} > \mathcal{C}_{D_2}$ and
therefore different effective thresholds.

\begin{predbox}
\textbf{CRF prediction}: a device with more measurement history
causes wave function collapse faster.

\[
  \tau = \tau_0 \cdot e^{-\kappa \cdot N_{\mathrm{eff}}}
  \qquad
  \kappa = \frac{k_B \ln(2)}{H_{\mathrm{material}}}
  \qquad
  N_{\mathrm{eff}} = \sum_i \Delta D_{\mathrm{KL},i}
\]

All parameters are fundamental constants ($k_B$, $\ln 2$) or
measurable physical quantities ($H_{\mathrm{material}}$,
$N_{\mathrm{eff}}$). Zero free parameters.

Standard QM: $\tau_1 = \tau_2$ (history-independent).\\
CRF: $\tau_1 < \tau_2$ when $N_{\mathrm{eff},1} > N_{\mathrm{eff},2}$.
\end{predbox}

\subsection{Distinction from Decoherence}

Decoherence theory predicts that two detectors with identical
\textit{current} physical state are equivalent, regardless of
history. CRF distinguishes history from current state:
$\mathcal{C}$ is encoded in $\Pr(s|\mathcal{C}_D)$ through prior
$\mathcal{R}$-events and persists beyond thermodynamic resets.
This is the empirical departure point.

\subsection{Experimental Protocol}

\textbf{Detector}: SNSPD (Superconducting Nanowire Single-Photon
Detector), NbN or WSi nanowire.

\textbf{Source}: SPDC entangled photon pairs; signal photons toward
detector; idler photons for coincidence timing.

\textbf{Setup}: $D_1$ (seasoned): $N_1 \sim 10^8$ prior
measurements at the experimental wavelength.
$D_2$ (fresh): $N_2 \ll N_1$, identical specifications.
Both thermally equilibrated before each run.

\textbf{Measurement}: collapse time $\tau$ (photon arrival to
detection event, resolution $<1$\,ps) for $\sim 10^6$ events per
block, alternating $D_1$ and $D_2$ in randomized blocks.
Blind analysis: detector identity concealed until after
distributions are computed.

\textbf{Signal}: QM: $\tau_1 = \tau_2$.
CRF: $\tau_1 < \tau_2$ at $>5\sigma$.

\textbf{Falsification}: a null result at sufficient statistical
power requires revision of the $\mathcal{C}$-accumulation
mechanism, though not the core derivation of $\theta$ from
$D_{\mathrm{KL}}$.

%% ── 6 ───────────────────────────────────────────────────────────────────
\section{Summary of Results}

\begin{center}
\renewcommand{\arraystretch}{1.6}
\begin{tabular}{p{5.5cm} p{2.2cm} p{4.5cm}}
\toprule
\textbf{Result} & \textbf{Status} & \textbf{Note} \\
\midrule
$\delta$ as Indefinable Primitive
  & Declared & Foundation of all derivations \\
$\theta = 1 - e^{-D_{\mathrm{KL}}/D_0}$
  & Derived & From boundary conditions on $f$ \\
$\alpha = \ln(2)$
  & Derived & $\delta$-bit equivalence \\
$\kappa = k_B\ln(2)/H_{\mathrm{material}}$
  & Derived & No free parameters \\
Born Rule $= |\langle s|\psi\rangle|^2$
  & Theorem & Via Gleason + SO(3) \\
Schwarzschild metric
  & Derived & From non-linear $\mathcal{R}$-feedback \\
Horizon $= \mathcal{R} \to 0$
  & Derived & Universe stops deciding \\
$G = \ln(2)\cdot c^2/D_0$
  & Derived & No external gravity parameter \\
$G(t) \propto 1/P(t)$
  & Hypothesis & Testable; consistent with DESI 2024 \\
P0 prediction $\tau_1 < \tau_2$
  & Testable & SNSPD protocol ready \\
\bottomrule
\end{tabular}
\end{center}

%% ── 7 ───────────────────────────────────────────────────────────────────
\section{Discussion}

\subsection{On the Measurement Problem}

CRF dissolves the measurement problem at its source. An observer is
any system with sufficient accumulated $\mathcal{C}$ to lower the
local $\mathcal{R}$-threshold. The Born Rule is the statistics of
minimal-$\mathcal{C}$ Resolution. As $\mathcal{C}$ accumulates,
statistics depart from Born Rule --- which P0 tests.

\subsection{On the Gravitational Constant}

$G = \ln(2)\cdot c^2/D_0$ connects gravity to information geometry.
The smallness of $G$ (in SI units) reflects the largeness of $D_0$
--- the cosmos has accumulated enormous $\delta$-activity,
making $P$ large and gravity weak. A younger universe would have
larger $G$. The prediction $G(t) \propto 1/P(t)$ is testable
from precision cosmology over long baselines.

\subsection{On the Framework}

CRF does not begin with physics. It begins with the question of
what has to be true for anything to be different from anything else.

The prediction in Section 5 is not the point. It is a consequence
--- one among many that the framework generates. Its value is not
that it confirms CRF if true. Its value is that it exposes CRF to
the possibility of being wrong.

That exposure is what distinguishes a framework from a belief.

\bigskip
\begin{center}
\itshape
$\delta$ does not stop anywhere.\\
It does not start anywhere either.\\
It simply operates --- and everything follows.
\end{center}

\end{document}
